\section{A square-root angular law and quantitative rank escape}
\label{sec:analytic-angular-heat-rate}

The following argument uses A to retain the full real zero product
of $H_0$, with multiplicities. It strengthens NN-A020's fixed-sector
argument quantitatively. Write $x_j>0$ for the positive zeros of
$H_0$, including multiplicity, and
\[
 H_0(z)=H_0(0)\prod_j(1-z^2/x_j^2),\qquad
 N(X)=\#\{j:x_j\le X\}.
\]
The product has no residual exponential factor: as in NN-A007,
$H_0(i\sqrt w)/H_0(0)$ has order at most $1/2<1$, so its genus-zero
Hadamard factorization has constant zero-free factor, fixed at one.
In particular $\sum_jx_j^{-2}<\infty$. The classical zero-count
bound supplies an effectively computable constant $C_0$ with
\[
                       N(X)\le C_0 X\log(X+3)\qquad(X>0).
\]
For the explicit primary bound one can use
Hasanalizade--Shen--Wong, Corollary~1.2, equation~(1.5),
printed p.~2, with $x_j=2\gamma_j$, as already checked in
NN-A014. A compact initial interval is covered by an effective
bound and a certified zero-free neighbourhood of zero.

\begin{lemma}[A global canonical-product remainder]
\label{lem:analytic-product-remainder}
For a nonreal $z$, set
\[
 L(z)=\frac{H_0'}{H_0}(z),\qquad
 S(z)=\sum_j\left(\frac1{|z-x_j|^2}+\frac1{|z+x_j|^2}\right)>0.
\]
The entire function
\[
 R_z(h)=\frac{H_0(z+h)}{H_0(z)e^{L(z)h}}
\]
satisfies, for every $h\in\mathbb C$,
\begin{align}
 |R_z(h)|&\le \exp\bigl(S(z)|h|^2/2\bigr),
                  \label{eq:analytic-R-global}\\
 |R_z(h)-1|&\le
     8S(z)|h|^2\exp\bigl(2S(z)|h|^2\bigr).
                  \label{eq:analytic-R-difference}
\end{align}
\end{lemma}

\begin{proof}
The paired product and logarithmic derivative give the locally
uniformly convergent canonical product
\[
 R_z(h)=\prod_{\xi\in\{x_j,-x_j\}}
           \left(1+\frac h{z-\xi}\right)
                       \exp\left(-\frac h{z-\xi}\right).
\]
The sum of $|z-\xi|^{-2}$ is finite; pairing makes the relation
with $L(z)$ explicit. No unrecorded exponential term remains.
For every complex $v$,
\[
 |1+v|^2\le\exp(2\operatorname{Re}v+|v|^2),
 \qquad |(1+v)e^{-v}|\le e^{|v|^2/2}.
\]
Multiply this inequality over finite products and take their
locally uniform limit to obtain \eqref{eq:analytic-R-global}.

Also $R_z(0)=1$ and $R_z'(0)=0$. Put $S=S(z)$ and
$Q(v)=R_z(v/\sqrt S)$. On $|v|=1$, $|Q(v)|\le e^{1/2}$.
Cauchy's coefficient estimate and the vanishing linear
coefficient give, for $|v|\le1/2$,
\[
 |Q(v)-1|
       \le e^{1/2}\sum_{k\ge2}|v|^k
       \le2e^{1/2}|v|^2<4|v|^2.
\]
For $|v|\ge1/2$, the first bound instead gives
$|Q(v)-1|\le2e^{|v|^2/2}\le8|v|^2e^{2|v|^2}$.
The same latter bound covers $|v|\le1/2$, proving
\eqref{eq:analytic-R-difference}.
\end{proof}

\begin{lemma}[A numerical zero-free test for the Gaussian convolution]
\label{lem:analytic-gaussian-nonzero}
For $a>0$ and nonreal $z$, put
\[
                A_z=aS(z),\qquad B_z=a^2|L(z)|^2S(z).
\]
If $A_z\le2^{-10}$ and $B_z\le2^{-10}$, then $H_{-a}(z)\ne0$.
More precisely,
\begin{equation}\label{eq:analytic-relative-Gaussian-bound}
 \left|\frac{H_{-a}(z)}{H_0(z)e^{aL(z)^2}}-1\right|
 \le16e^{16B_z}
 \left\{\frac{4B_z}{\sqrt{1-16A_z}}+
             \frac{2A_z}{(1-16A_z)^{3/2}}\right\}
 \le\frac7{16}<1.
\end{equation}
\end{lemma}

\begin{proof}
The theta integral gives the exact convolution
\[
 H_{-a}(z)=\frac1{\sqrt{4\pi a}}
             \int_{\mathbb R}e^{-u^2/(4a)}H_0(z+u)\,du.
\]
Let $p=2aL(z)$ and translate the integration line to
$\mathbb R+p$. The integrand is entire. On the vertical ends,
the Gaussian quadratic decay dominates the order-one bound
$|H_0(w)|\le\exp(C|w|\log(|w|+3))$ for fixed $z,p,a$.
Thus the end integrals tend to zero. A real part of $p$ only
reparametrizes the translated horizontal line. There is no
logarithm to continue and no obstruction from zeros crossed
in this deformation: $R_z$ itself is entire.
Completing the square gives
\[
 \frac{H_{-a}(z)}{H_0(z)e^{aL(z)^2}}
          =\mathbb E R_z(p+U),\qquad U\sim N(0,2a).
\]
With $S=S(z)$, \eqref{eq:analytic-R-difference} and
$|p+U|^2\le2|p|^2+2U^2$ show that the error is at most
\[
 16S e^{4S|p|^2}
       \mathbb E\bigl[(|p|^2+U^2)e^{4SU^2}\bigr].
\]
The two elementary Gaussian integrals are
\[
 \mathbb E e^{4SU^2}=(1-16aS)^{-1/2},\qquad
 \mathbb E U^2e^{4SU^2}
                       =2a(1-16aS)^{-3/2}.
\]
Since $S|p|^2=4B_z$, these give the first inequality in
\eqref{eq:analytic-relative-Gaussian-bound}.
For $A_z,B_z\le2^{-10}$, bound $e^{16B_z}<2$,
$(1-16A_z)^{-1/2}<2$, and
$(1-16A_z)^{-3/2}<3$. The right side is then at most
$256B_z+192A_z\le448/1024=7/16$.
The denominator is nonzero under A, and this strict relative
bound proves the assertion.
\end{proof}

\begin{lemma}[Uniform cone bounds with a vanishing radial tail]
\label{lem:analytic-cone-logderivative}
There is an effectively computable $C\ge1$ such that, whenever
$0<q\le1$, $r=|z|$ and $|\operatorname{Im}z|\ge qr$,
\begin{align}
 S(z)&\le Cq^{-2}\frac{\log(r+3)}{r+1},
                 \label{eq:analytic-S-cone}\\
 |L(z)|&\le Cq^{-1}\log(r+3).
                 \label{eq:analytic-L-cone}
\end{align}
Consequently
\[
 |L(z)|^2S(z)
       \le C^3q^{-4}\frac{\log^3(r+3)}{r+1}.
\]
\end{lemma}

\begin{proof}
By evenness and conjugation we may assume
$\operatorname{Re}z\ge0$. Choose a fixed positive $r_0$ with
$2r_0<x_1$. For $r\le r_0$, the product gives bounded $S(z)$
and $L(z)=O(r)$, which are absorbed by the claimed constants.
For $r\ge r_0$ split the positive zeros into
$x_j<r/2$, $r/2\le x_j\le2r$, and $x_j>2r$.

In the first range, $|z\pm x_j|\ge r/2$. The contribution to
$S$ is at most $8N(r/2)/r^2\ll\log(r+3)/r$.
Using the paired logarithmic derivative
\[
                  L(z)=-2z\sum_j\frac1{x_j^2-z^2},
\]
the same range contributes $O(N(r/2)/r)=O(\log(r+3))$.
In the middle range, $|z-x_j|\ge qr$ and, since
$\operatorname{Re}z\ge0$, $|z+x_j|\ge r$.
Its contributions are at most
$2q^{-2}N(2r)/r^2$ to $S$ and
$2q^{-1}N(2r)/r$ to $|L|$.
Finally partial summation of the count bound gives
\[
             \sum_{x_j>2r}x_j^{-2}
                          \ll\frac{\log(r+3)}r.
\]
This last range contributes at most eight times that sum to
$S$, and $O(r)$ times that sum to $|L|$.
Combine the bounds and compare $r$ with $r+1$ for $r\ge r_0$.

All constants can be made effective from the stated zero-count
bound, the computable sum $\sum_jx_j^{-2}=F_0'(0)$, and a
certified zero-free disc at zero. To obtain the last disc one
can bound $|H_0(z)-H_0(0)|$ by $|z|$ times the computable theta
integral bound for $H_0'$ on $|z|\le1$, and use $H_0(0)>0$.
No simplicity or root-separation estimate is used.
\end{proof}

\begin{theorem}[NN-A022: angular rate and detecting-order lower bounds]
Assume A and define
\[
 \alpha(-a)=
 \sup_{\substack{H_{-a}(z)=0\\\operatorname{Re}z>0}}|\arg z|,
                          \qquad a>0.
\]
There are effectively computable constants $a_0,c,c_U,C>0$ such
that for $0<a\le a_0$,
\begin{equation}\label{eq:analytic-root-angle-rate}
                         \alpha(-a)\le C\sqrt a.
\end{equation}
More precisely the high zeros have a uniformly vanishing
scaled contribution:
\begin{equation}\label{eq:analytic-angular-tail-small}
 \lim_{R\to\infty}
 \sup_{0<a\le1}\ 
 \sup_{\substack{H_{-a}(z)=0,\ \operatorname{Re}z>0\\|z|\ge R}}
                       \frac{|\arg z|}{\sqrt a}=0,
\end{equation}
where the supremum over an empty set is zero.

For every integer $r\ge1$, all allowed coefficient Toeplitz
minors of $F_t$ of order at most $r$, at arbitrary indices,
are strictly positive throughout
\[
                         -\frac c{(r+1)^2}\le t\le0.
\]
For every $d\ge1$, $\chi_{t,d}$ from NN-A021 is of positive
type, with every polynomial Schur weight strictly positive,
throughout
\[
                         -\frac {c_U}{(d+1)^2}\le t\le0.
\]
Thus the least negative-minor order and the least failing
unitary rank satisfy
\[
                 m_*(-a)\gg a^{-1/2},
                 \qquad d_*(-a)\gg a^{-1/2}\qquad(a\downarrow0).
\]
The least detecting index window of NN-A014 also satisfies
$R_*(-a)\ge m_*(-a)-1\gg a^{-1/2}$, since a minor in a window
of size $R+1$ has order at most $R+1$.
Rodgers--Tao is needed only to know that these detecting
quantities are finite at every negative time.
\end{theorem}

\begin{proof}
The functions
$\log(r+3)/(r+1)$ and $\log^3(r+3)/(r+1)$ are bounded on
$[0,\infty)$ and tend to zero at infinity. The last lemma
therefore gives absolute effective constants $C_1,C_2$ with
\[
                  S(z)\le C_1q^{-2},\qquad
                  |L(z)|^2S(z)\le C_2q^{-4}
\]
on the stated cones. Choose an effective $K\ge1$ such that
$C_1/K^2\le2^{-10}$ and $C_2/K^4\le2^{-10}$.
For $q=K\sqrt a\le1/2$, Lemma~\ref{lem:analytic-gaussian-nonzero}
excludes every zero with $|\operatorname{Im}z|\ge q|z|$.
There are no imaginary-axis zeros, since $H_{-a}(iy)>0$.
Hence
\[
                 \alpha(-a)\le\arcsin(K\sqrt a)\le2K\sqrt a
                    \qquad(0<a\le1/(4K^2)).
\]
This proves \eqref{eq:analytic-root-angle-rate}.

For the sharper tail statement, fix $0<\varepsilon<1/2$.
In the cone with $q=\varepsilon\sqrt a$, $0<a\le1$, the
two parameters in the Gaussian nonvanishing test obey
\[
 A_z\le\frac C{\varepsilon^2}
               \frac{\log(r+3)}{r+1},\qquad
 B_z\le\frac{C^3}{\varepsilon^4}
               \frac{\log^3(r+3)}{r+1}.
\]
Both are at most $2^{-10}$ for $r\ge R_\varepsilon$,
uniformly in $a$. A zero of that size must therefore have
$|\arg z|\le\arcsin(\varepsilon\sqrt a)
\le2\varepsilon\sqrt a$. Let $\varepsilon$ decrease to zero
to obtain \eqref{eq:analytic-angular-tail-small}.

The quadratic factor criterion in NN-A020 applies whenever
$2\alpha(t)<\pi/(r+1)$. The displayed angular bound supplies
this inequality on $-c/(r+1)^2\le t\le0$ after choosing a
fixed smaller positive effective $c$. The infinite positive
particle paths give strictness for every allowed minor.
For example, with the preceding $K$, taking
$c\le\pi^2/(64K^2)$ gives a factor-two margin in the angle
inequality and lies inside the small-time range.

For the unitary statement retain the effective constants
$b<1$ and $\eta_0>0$ in the proof of NN-A021, so
$|\beta_j(t)|<b$ when $-\eta_0\le t\le0$. Then
\[
          |\arg\widetilde\beta_j(t)|
               \le\frac{2\alpha(t)}{1-b}
               \le\frac{4K\sqrt{|t|}}{1-b}.
\]
Choose $c_U>0$ with $c_U\le\eta_0$ and
$c_U\le\pi^2(1-b)^2/(64K^2)$. The same quadratic-factor
argument and the exact finite-rank dictionary of NN-A021
give the conclusion for every $d$. Taking the largest
integers $r,d$ permitted by these intervals gives the two
lower bounds for the detecting order and rank.
\end{proof}

\begin{corollary}[One explicit numerical certificate]
The constant in the $PF_r$ interval may be chosen as
$c=2^{-36}$. Thus, under A, every allowed minor of order at
most $r$ is strictly positive for
\[
                       -\frac{2^{-36}}{(r+1)^2}\le t\le0.
\]
For $0<a\le2^{-38}$ this implies
\[
 m_*(-a)\ge
       \left\lfloor\frac{2^{-18}}{\sqrt a}\right\rfloor,
 \qquad
 R_*(-a)\ge
       \left\lfloor\frac{2^{-18}}{\sqrt a}\right\rfloor-1.
\]
The numerical input is the existing certified inequality
$A_0<1$ from NN-A014; no new floating-point experiment is used.
\end{corollary}

\begin{proof}
The coefficient sum from NN-A014 obeys
\[
 A_0=\sum_{k\ge0}(2k+2)(2k+1)c_{k+1}(0)
                              \ge2c_1(0)=2p_1(0).
\]
Its saved Arb certificate gives $0<p_1(0)<1/2$. For real
$|x|\le1$, the elementary inequality $\cos v\ge1-v^2/2$
therefore implies
\[
 H_0(x)\ge H_0(0)(1-p_1(0)x^2)>0.
\]
In particular $x_1>1$, and $\sum_jx_j^{-2}=p_1(0)<1/2$.

Write $N_\zeta(T)$ for the positive-ordinate zeta count.
The exact source bound, for $T\ge e$, is
\[
 \left|N_\zeta(T)-\frac{T}{2\pi}
                      \log\frac{T}{2\pi e}\right|
       \le .1038\log T+.2573\log\log T+9.3675.
\]
This is Hasanalizade--Shen--Wong, Corollary~1.2,
equation~(1.5), printed p.~2. At $T=e$ the main term is
negative, so $N_\zeta(e)<10$. For $T\ge e$, the main term
is at most $T\log(T+3)$, and the error is less than
\[
 .362\log T+9.368
       \le .362T+4T
       \le4.362T\log(T+3).
\]
Here $\log\log T\le\log T$, $e>5/2$, and
$\log(T+3)>1$ suffice. Hence
$N_\zeta(T)\le6T\log(T+3)$. Together with $x_1>1$ and
$N(X)=N_\zeta(X/2)$, these bounds give the deliberately loose
but uniform inequality
\begin{equation}\label{eq:analytic-explicit-count-100}
                         N(X)\le100X\log(X+3)\qquad(X>0).
\end{equation}
Indeed for $X<1$ the left side is zero, for
$1\le X<2e$ it is less than ten, and for $X\ge2e$
it is at most $3X\log(X+3)$.

We now make every constant in the cone lemma explicit.
For $r\ge1/2$, the range $x_j<r/2$ contributes at most
\[
       400\frac{\log(r+3)}r\quad\hbox{to }S,
       \qquad \frac{400}3\log(r+3)\quad\hbox{to }|L|.
\]
The middle range contributes at most
\[
       800q^{-2}\frac{\log(r+3)}r\quad\hbox{to }S,
       \qquad800q^{-1}\log(r+3)\quad\hbox{to }|L|,
\]
using $\log(2r+3)\le2\log(r+3)$.
For the last range, partial summation and
$\log(2r+3)>1$ give
\begin{align*}
 \sum_{x_j>2r}x_j^{-2}
 &\le200\int_{2r}^\infty\frac{\log(u+3)}{u^2}\,du\\
 &\le200\,\frac{\log(2r+3)+1}{2r}
 \le400\,\frac{\log(r+3)}r.
\end{align*}
It contributes at most
$3200\log(r+3)/r$ to $S$ and
$(3200/3)\log(r+3)$ to $|L|$.
Thus
\[
 S\le4400q^{-2}\frac{\log(r+3)}r
       \le13200q^{-2}\frac{\log(r+3)}{r+1},
       \qquad |L|\le2000q^{-1}\log(r+3).
\]
For $r\le1/2$, $x_1>1$ and $p_1(0)<1/2$ give
$S<4$ and $|L|\le4r/3\le2/3$, which satisfy the same
bounds after taking $C=2^{14}$ in the cone lemma.

For every $r\ge0$,
\[
 \frac{\log(r+3)}{r+1}\le2,\qquad
 \frac{\log^3(r+3)}{r+1}\le16.
\]
For the second inequality put $y=r+3\ge3$, note
$r+1=y-2\ge y/3$, and maximize $(\log y)^3/y$:
its maximum is $27/e^3$, so the ratio is at most
$81/e^3<81/8<16$. The first follows from
$\log(r+3)\le r+2$.
Consequently one can use $C_1=2^{15}$ and $C_2=2^{46}$
in the proof of the theorem. With $K=2^{16}$,
\[
             C_1/K^2=2^{-17},\qquad C_2/K^4=2^{-18},
\]
both smaller than the required $2^{-10}$. Thus
\[
 \alpha(-a)\le2^{17}\sqrt a\qquad(0<a\le2^{-34}).
\]
If $a\le2^{-36}/(r+1)^2$, then
$2\alpha(-a)\le1/(r+1)<\pi/(r+1)$.
The quadratic factor criterion proves the explicit interval.
Finally take
$r=\lfloor2^{-18}/\sqrt a\rfloor-1\ge1$.
All minors of order at most this $r$ are positive, so
$m_*(-a)\ge r+1$, and $R_*(-a)\ge m_*(-a)-1$.
\end{proof}

\begin{corollary}[The limiting angle detects endpoint multiplicities]
For each distinct positive zero $x$ of $H_0$, let $m_x$ be its
multiplicity, and let $h_m$ be the largest absolute value of a
root of the probabilists' Hermite polynomial $\operatorname{He}_m$;
in particular $h_1=0$. Under A,
\begin{equation}\label{eq:analytic-endpoint-angular-limit}
       \lim_{a\downarrow0}\frac{\alpha(-a)}{\sqrt a}
          =\kappa_0,\qquad
       \kappa_0=\sup_{H_0(x)=0,\ x>0}\frac{\sqrt2\,h_{m_x}}x.
\end{equation}
The supremum is finite and is attained if it is positive.
In particular all nontrivial zeta zeros are simple if and only
if $\alpha(-a)=o(\sqrt a)$. Under that additional simplicity
condition the two detection bounds improve to
\[
                     \sqrt a\,m_*(-a)\longrightarrow\infty,
              \qquad \sqrt a\,d_*(-a)\longrightarrow\infty.
\]
This does not determine the height of the first nonreal zero.
\end{corollary}

\begin{proof}
The Hermite recurrence gives its roots as the eigenvalues of
the real symmetric tridiagonal matrix with off-diagonal
entries $\sqrt1,\ldots,\sqrt{m-1}$. Gershgorin's bound yields
$h_m\le2\sqrt m$. Since
$m_x\le N(2x)\ll x\log(x+3)$, it follows that
$h_{m_x}/x\to0$ as $x\to\infty$. There are finitely many
distinct zeros in each compact set, so $\kappa_0$ is finite
and any positive supremum is a maximum.

At each fixed positive zero $x$, NN-A006 gives the complete
local splitting for negative time:
\[
             z_j(-a)=x+i\sqrt{2a}\,h_{m_x,j}+O(a).
\]
Consequently the largest angle of these local branches,
divided by $\sqrt a$, tends to $\sqrt2\,h_{m_x}/x$.
This also covers a simple zero, whose branch remains real
for small real time.

Choose a radius $R$ whose boundary contains no zero of $H_0$.
Uniform compact convergence, local root splitting and
Rouch\'e's theorem account for every zero in $|z|\le R$.
There are finitely many endpoint germs there, so the scaled
maximum of their angles tends to the maximum of their
Hermite constants. The remaining zeros contribute arbitrarily
little, uniformly for $0<a\le1$, by
\eqref{eq:analytic-angular-tail-small}. Combining the finite
maximum with this vanishing tail proves
\eqref{eq:analytic-endpoint-angular-limit}.
Since $h_m>0$ precisely when $m\ge2$, $\kappa_0=0$ is
equivalent to simplicity of every endpoint zero.

If $\alpha(-a)/\sqrt a\to0$, then for any prescribed $M>0$
the quadratic angle criterion applies for all orders
$r\le M/\sqrt a$ once $a$ is sufficiently small. The
M\"obius angle bound from NN-A021 gives the same assertion
at unitary ranks $d\le M/\sqrt a$. Since $M$ is arbitrary,
the two strengthened lower bounds follow.
\end{proof}

\paragraph{Exact geometric and representation meaning.}
The multiplicity $m_x$ is the length of the completed local
zero algebra. The parameter change $t=-u^2/2$ splits its
initial heat polynomial into the distinct Hermite directions;
the positive real direction $x$ converts their imaginary
normal displacement into the coefficient
$\sqrt2\,h_{m_x}/x$. The global theorem shows that these
finite local Artin lengths determine the limiting maximal
angle, because the high-radius contribution vanishes uniformly.
The lower bounds concern every skew Schur test of bounded
width and every polynomial irreducible at bounded unitary
rank, even when its degree is unrestricted. Their
exterior-power dictionary is SGA~6, Expos\'e~V,
Example~2.2(b) and Lemma~2.2.1, with the exact finite-rank
equivalence proved in NN-A021. The bounds are for these
specified certificate families; no general complexity lower
bound, pure arithmetic motive, or full collision-height law
is asserted.
